\documentclass[11pt]{article}
\usepackage{amsmath, amssymb, graphicx, geometry, hyperref}
\geometry{margin=1in}

\title{\textbf{DRIFT: A Phase-Space Monitoring Framework for Earth Rotation Dynamics}}
\author{Craig Stone}
\date{\today}

\begin{document}

\maketitle

\begin{abstract}
We present the DRIFT (Directional Rotational Inference and Forecasting Toolkit) dashboard, a phase-space–based monitoring system for Earth’s polar motion. Departing from conventional additive forcing models, DRIFT adopts a geometry-first approach, identifying invariant structure directly from observational data. The system reveals a low-dimensional organization of polar motion, characterized by a dominant axis and a two-state regime structure. By embedding this structure into a real-time diagnostic and forecasting framework, DRIFT enables early detection of rotational transitions through phase-space deviations, lag-response structure, and alignment diagnostics. This work formalizes both the theoretical basis and implementation of DRIFT, and provides a panel-by-panel interpretation of its analytical components.
\end{abstract}

\section{Introduction}

Polar motion is traditionally modeled as the superposition of atmospheric, oceanic, and cryospheric forcings. This approach presupposes causation and distributes explanatory weight across multiple partially constrained inputs.

In contrast, this work inverts the problem. Rather than assuming causes, we first identify the geometric and dynamical constraints required by the observed data. Over the modern observational record, polar motion exhibits:

\begin{itemize}
\item Confinement to a low-dimensional manifold
\item Alignment along a dominant Earth-fixed axis
\item Intermittent switching between two dynamical states
\end{itemize}

The DRIFT framework operationalizes these findings into a monitoring and forecasting system.

\section{Conceptual Framework}

\subsection{Phase-Space Representation}

Polar motion is mapped into a phase-space representation:

\begin{equation}
\theta(t) = \arctan\left(\frac{y_p(t)}{x_p(t)}\right), \quad \omega(t) = \frac{d\theta}{dt}
\end{equation}

This transformation reveals cyclic structure and regime transitions not apparent in Cartesian coordinates.

\subsection{Two-State Dynamics}

The system exhibits two persistent modes:

\begin{itemize}
\item \textbf{Stable Mode:} Regular cyclic motion with tight phase loops
\item \textbf{Transition Mode:} Distorted trajectories, phase slips, and lag restructuring
\end{itemize}

\subsection{Geometry-First Inference Hierarchy}

\begin{enumerate}
\item \textbf{Geometric invariants} (axes, confinement)
\item \textbf{Dynamical organization} (cycles, states)
\item \textbf{Directional features} (alignment, drift)
\end{enumerate}

DRIFT is constructed to respect this hierarchy.

\section{System Architecture}

The dashboard integrates:

\begin{itemize}
\item Time-series polar motion data
\item Derived phase-space variables
\item Geomagnetic proxies
\item Lag-response modeling
\item Real-time state inference
\end{itemize}

Each panel represents a projection of the system into a diagnostic subspace.

\section{Panel-by-Panel Description}

\subsection{1. Transition Forecast}

This panel estimates the probability of an upcoming regime transition.

\begin{itemize}
\item The probability curve represents a forward convolution of lag-response structure.
\item Peak time indicates expected transition horizon.
\item Cumulative probability tracks integrated risk.
\end{itemize}

Interpretation:
High early probability indicates imminent transition; delayed peaks suggest latent instability.

\subsection{2. 3D Vector View}

Displays:

\begin{itemize}
\item Drift vector
\item Principal axes ($e_1, e_2, e_3$)
\item Geomagnetic dipole orientation
\end{itemize}

Interpretation:
Small angular separations indicate coupling; divergence suggests decoupling or reorientation.

\subsection{3. Polar Motion (XP, YP)}

Raw observational data with turning points marked.

\begin{itemize}
\item Arc structure reveals dominant direction
\item Turning points mark phase inflections
\end{itemize}

Interpretation:
Tight arcs indicate stability; distortions precede transitions.

\subsection{4. Drift Direction}

Tracks longitude of drift axis over time.

\begin{itemize}
\item Persistent slope = steady migration
\item Plateaus = regime locking
\end{itemize}

Interpretation:
Sharp changes indicate structural reorganization.

\subsection{5. Angle Diagnostics}

Tracks relative angular relationships:

\begin{itemize}
\item $\theta_3$: orientation vs principal axis
\item $\theta_{12}$: in-plane alignment
\end{itemize}

Interpretation:
Spikes and clustering indicate instability or reorientation.

\subsection{6. Alignment: Drift vs Geomagnetic}

Compares drift axis to geomagnetic activity indices.

\begin{itemize}
\item Lower values = stronger alignment
\item Overlays with Kp/ap indices
\end{itemize}

Interpretation:
Coincident shifts suggest coupling between rotational and geomagnetic systems.

\subsection{7. Phase Portrait ($\theta$ vs $\omega$)}

Core diagnostic.

\begin{itemize}
\item Tight loops = stable periodic motion
\item Distorted loops = transition regime
\end{itemize}

Interpretation:
Topology change in loops signals regime shift.

\subsection{8. Phase Diagnostics ($\theta(t), \omega(t)$)}

Time evolution of phase variables.

\begin{itemize}
\item Phase clustering
\item Angular velocity modulation
\end{itemize}

Interpretation:
Slowdowns and irregularities precede transitions.

\subsection{9. Orthogonal Deviation Ratio $R(t)$}

Measures anisotropy:

\begin{equation}
R(t) = \frac{\min(\lambda_i)}{\max(\lambda_i)}
\end{equation}

Interpretation:
\begin{itemize}
\item Low $R(t)$: elongated, directional state
\item High $R(t)$: isotropic, unstable state
\end{itemize}

\subsection{10. Overlay Plot}

Normalized signals for comparison.

Interpretation:
Lead-lag relationships become visually apparent.

\subsection{11. Lag Model}

Impulse-response estimate.

\begin{itemize}
\item Shows delayed system response after turning points
\item Baseline vs observed response
\end{itemize}

Interpretation:
Positive excursions indicate structured response, not noise.

\subsection{12. Conditional Lag Response}

Phase-conditioned lag structure.

\begin{itemize}
\item Heatmap reveals phase-dependent response
\item Lower panel shows slice through phase bin
\end{itemize}

Interpretation:
Confirms that system response depends on phase state, not just event occurrence.

\section{Forecasting Framework}

DRIFT forecasting emerges from:

\begin{itemize}
\item Detection of phase instability
\item Lag-response amplification
\item Alignment degradation
\end{itemize}

Transition probability is computed as:

\begin{equation}
P(t+\tau) = \int R_{\text{lag}}(\tau) \cdot I_{\text{instability}}(t) \, d\tau
\end{equation}

\section{Discussion}

DRIFT reframes polar motion as a constrained dynamical system rather than a sum of forcings. The emergence of:

\begin{itemize}
\item Earth-fixed geometric structure
\item Two-state dynamics
\item Phase-dependent response
\end{itemize}

suggests the presence of an underlying weakly coupled rotational system.

\section{Conclusion}

The DRIFT dashboard transforms polar motion analysis into a real-time diagnostic and forecasting tool. By grounding inference in geometry and phase-space structure, it provides a robust framework for identifying and anticipating rotational regime shifts.

\section*{Future Work}

\begin{itemize}
\item Integration with geomagnetic field models
\item Null-model validation using surrogate data
\item Extension to full inertial tensor estimation
\end{itemize}

\end{document}